(1) 相速度为 \(v_p=\omega(k_0)/k_0=c+\alpha k_0^2\)，群速度 \(v_g=\omega'(k_0)=c+3\alpha k_0^2\)；二阶系数 \(\omega_2=\omega''(k_0)=6\alpha k_0\)。三者承担不同角色：相位、包络中心与包络展宽。
(2) 写 \(k=k_0+\kappa\)，则 \(\omega(k)\approx\omega_0+v_g\kappa+\frac12\omega_2\kappa^2\)。傅里叶积分中线性相位为 \(\kappa(x-v_gt)\)，因此高斯包络中心位于 \(x_c=v_gt\)，总载波相位为 \(k_0x-\omega_0t\)；二阶项不会改变本近似下的中心速度。
(3) 初始包络的波数权重正比于 \(e^{-\sigma_0^2\kappa^2}\)。传播后高斯积分的复宽度参数为 \(D=\sigma_0^2+i\omega_2t/2\)，于是强度仍为高斯，其标准差 \(\sigma_x(t)=\sigma_0\sqrt{1+(\omega_2t/(2\sigma_0^2))^2}\)。自然展宽时间 \(t_D=2\sigma_0^2/|\omega_2|\)，二次相位曲率的符号随 \(\omega_2t\) 反转。
(4) \(\alpha\to0\) 时 \(\omega_2\to0\)，高斯宽度保持 \(\sigma_0\)，仅以 \(c\) 平移。为忽略三阶余项需谱宽 \(\Delta k\sim1/\sigma_0\ll k_0\) 且传播相位误差 \(|\omega'''(k_0)|t(\Delta k)^3/6\ll1\)；时间过长时即使初态窄带，也不能只依赖二阶高斯近似。
